% ======================================================================
%  PART 11 -- Linear Time-Invariant (LTI) Filters
% ======================================================================
\section{Linear Time-Invariant (LTI) Filters}
\label{sec:lti}

A \emph{(digital) filter} is a map that turns one sequence into another,
$\{Y_t\}=\filt[\{X_t\}]$, with index set $\Ints$. The filters that matter for
time series are the \textbf{linear time-invariant} (LTI) ones: they treat all
time points alike and act linearly. Two ``test inputs'' completely characterise
such a filter --- a single impulse (giving the \emph{impulse response} $h$) and
a pure complex wave (giving the \emph{transfer function} $H$). Their power is
twofold: every LTI filter is a \emph{convolution} with $h$, and in the frequency
domain it acts \emph{diagonally}, multiplying the spectral density by
$\abs{H(f)}^2$. This single fact, $\sdf_Y(f)=\abs{H(f)}^2\sdf_X(f)$, is the
engine behind the closed-form spectral density of any ARMA process. Throughout
we allow time series to be \emph{complex-valued} and assume the (possibly
infinite) sums involved converge, the standard sufficient condition being
$h\in\ell^1$.

% ----------------------------------------------------------------------
\subsection{Definitions}

\begin{definition}{Complex-valued time series \& its autocovariance}{p11-complex}
A \textbf{complex-valued} time series can be written $X_t=U_t+iV_t$ with
$\{U_t\},\{V_t\}$ real-valued, and has mean
$\E[X_t]=\E[U_t]+i\,\E[V_t]$. Its \textbf{autocovariance} uses a
\emph{conjugate} on the second factor:
\[
\acvs^{(X)}_\tau=\Cov(X_{t+\tau},X_t)=\E\!\big[(X_{t+\tau}-\mu)\conj{(X_t-\mu)}\big].
\]
This reduces to the usual real ACVS when $X_t\in\Reals$. (Stationarity is
required jointly of the real and imaginary parts.) Throughout we drop the index
set: $\{X_t\}=\{X_t\}_{t\in\Ints}$.
\end{definition}

\begin{definition}{Digital filter \& output notation}{p11-filter}
A \textbf{digital filter} $\filt$ maps an input sequence to an output sequence,
$\{Y_t\}=\filt[\{X_t\}]$. To name a single output entry we put the time index
\emph{after} the bracket,
\[
Y_u=\filt[\{X_t\}]_u,\qquad u\in\Ints .
\]
The $t$ in $\{X_t\}$ is a \emph{dummy} variable; the $u\in\Ints$ after the
bracket has a definite meaning (the output time).
\end{definition}

\begin{definition}{Linear time-invariant (LTI) filter}{p11-lti}
A digital filter $\filt$ is \textbf{linear time-invariant} if for all sequences
$\{X_t\},\{Y_t\}$ and all $\alpha\in\Comp$:
\begin{enumerate}
  \item \textbf{Linearity:}
        $\displaystyle \filt[\{\alpha X_t+Y_t\}]=\alpha\,\filt[\{X_t\}]+\filt[\{Y_t\}]$;
  \item \textbf{Time invariance:}
        $\displaystyle \filt\big[\Bsh[\{X_t\}]\big]=\Bsh\big[\filt[\{X_t\}]\big]$,
        i.e.\ $\filt$ \emph{commutes with the backshift}, $\filt\Bsh=\Bsh\filt$.
\end{enumerate}
Time invariance for the single shift $\Bsh$ already implies it for every shift:
$\filt[\Bsh^u\{X_t\}]=\Bsh^u\filt[\{X_t\}]$ for all $u\in\Ints$. (The
continuous-time analogue is usually stated with $\Bsh^u$, but the weaker
one-step condition suffices.)
\end{definition}

\begin{definition}{Impulse sequence \& impulse response}{p11-impulse}
For $m\in\Ints$ the \textbf{impulse sequence} $\{\delta_{t,m}\}$ is
$\delta_{t,m}=1$ if $t=m$ and $0$ otherwise. The \textbf{impulse response} of an
LTI filter $\filt$ is the sequence $\{h_m\}$ defined by
\[
h_m=\filt[\{\delta_{t,-m}\}]_0,\qquad m\in\Ints,
\]
where the subscript $0$ is the \emph{output time index}. Informally $h_m$ is the
response, read off at time $0$, to a unit shock placed at time $-m$. A common
sufficient regularity condition is $h\in\ell^1$
(then $\norm{h}_\infty\le\norm{h}_1<\infty$).
\end{definition}

\begin{definition}{Complex wave \& transfer function}{p11-transfer}
For $f\in\Reals$ the \textbf{complex wave} is the sequence
$\{\xi_{t,f}\}$ with $\xi_{t,f}=e^{2\pi i f t}$. The \textbf{transfer function}
of an LTI filter $\filt$ that can be applied to these waves is
\[
H(f)=\filt[\{\xi_{t,f}\}]_0,\qquad f\in\Reals .
\]
It encodes the frequency-domain behaviour of $\filt$: how each pure frequency is
re-weighted in amplitude and shifted in phase.
\end{definition}

% ----------------------------------------------------------------------
\subsection{Theorems \& Propositions}

\begin{proposition}{Linear combination of LTI filters is LTI}{p11-lincomb}
If $\filt_1,\filt_2$ are LTI and $\alpha\in\Comp$, then
$\filt=\alpha\filt_1+\filt_2$ is LTI.\\[2pt]
\textit{Proof idea:} apply linearity of $\filt_1,\filt_2$ to
$\filt[\beta\{X_t\}+\{Y_t\}]$ and regroup; for time invariance use that each
$\filt_j$ commutes with $\Bsh$ and that $\Bsh$ is linear, so
$\filt[\Bsh\{X_t\}]=\Bsh[\filt\{X_t\}]$.\\
\textit{Why:} it makes the LTI filters a \emph{vector space} (over $\Comp$); in
particular any polynomial in $\Bsh$, e.g.\ $\Phi(\Bsh)=\sum_j\phi_j\Bsh^j$, is
LTI.
\end{proposition}

\begin{proposition}{Cascade of LTI filters is LTI}{p11-cascade}
If $\filt_1,\filt_2$ are LTI, then the \textbf{cascade}
$\filt=\filt_1\filt_2$, defined by
$\filt[\{X_t\}]=\filt_1\big[\filt_2[\{X_t\}]\big]$, is LTI.\\[2pt]
\textit{Proof idea:} linearity is the composition of two linear maps; time
invariance follows by pushing $\Bsh$ through both filters,
$\filt_1[\filt_2[\Bsh\{X_t\}]]=\filt_1[\Bsh[\filt_2\{X_t\}]]
=\Bsh[\filt_1[\filt_2\{X_t\}]]$.\\
\textit{Why:} filters can be \emph{chained} and the family is closed under
composition; combined with the previous proposition, all of $\{$polynomials and
power series in $\Bsh\}$ are LTI.
\end{proposition}

\begin{theorem}{LTI $\Leftrightarrow$ convolution with the impulse response}{p11-conv}
Under suitable convergence assumptions, an LTI filter with impulse response $h$
acts as a \textbf{convolution}:
\[
\filt[\{X_t\}]_u\ms\sum_{m\in\Ints}h_{u-m}X_m,\qquad u\in\Ints .
\]
Conversely, any filter defined by such a convolution (with, say, $h\in\ell^1$)
is LTI. The impulse response gives the convolution weights.\\[2pt]
\textit{Proof idea (deterministic $\ell^1$ version):}
write $X^{(N)}=\sum_{\abs m\le N}X_m\delta_{\cdot,m}$; by continuity in
$\ell^1$, $\filt[\{X_t\}]_u=\lim_N\filt[X^{(N)}]_u$. Linearity turns each finite
truncation into a finite sum, and time invariance gives
$\filt[\{\delta_{t,m}\}]_u=\filt[\{\delta_{t,m-u}\}]_0=h_{u-m}$; absolute
convergence ($h\in\ell^1$, $X\in\ell^1$) lets the partial sums pass to the full
series. The converse checks linearity and time invariance directly on the
convolution (re-indexing $m'=m-1$ for the shift).\\
\textit{Why:} this is the master representation --- it reduces \emph{any} LTI
filter to a single object, the weights $h$, and underlies every spectral
computation below.
\end{theorem}

\begin{theorem}{Complex waves are eigensequences; $H(f)$ is the eigenvalue}{p11-eigen}
Let $\filt$ be an LTI filter applicable to the complex waves $\{\xi_{t,f}\}$.
Then each wave is an \textbf{eigensequence}: for every $f\in\Reals$,
\[
\filt[\{\xi_{t,f}\}]=H(f)\,\{\xi_{t,f}\},
\]
i.e.\ $\filt$ only rescales amplitude and shifts phase while preserving the
frequency.\\[2pt]
\textit{Proof idea:} since
$\Bsh^{-u}[\{\xi_{t,f}\}]_\tau=\xi_{\tau+u,f}=e^{2\pi i u f}\xi_{\tau,f}$, we have
$\Bsh^{-u}[\{\xi_{t,f}\}]=e^{2\pi i u f}\{\xi_{t,f}\}$. Then for any $u$,
\[
\filt[\{\xi_{t,f}\}]_u
=\Bsh^{u}\!\big[\filt[\Bsh^{-u}\{\xi_{t,f}\}]\big]_u
=\filt\big[e^{2\pi i f u}\{\xi_{t,f}\}\big]_0
=e^{2\pi i f u}\filt[\{\xi_{t,f}\}]_0=\xi_{u,f}\,H(f),
\]
using time invariance, the shift identity, linearity and the definition of $H$.\\
\textit{Why:} it diagonalises every LTI filter in the frequency basis ---
filtering becomes \emph{pointwise multiplication} by $H(f)$, the source of all
clean frequency-domain formulas.
\end{theorem}

\begin{theorem}{Transfer function $=$ Fourier transform of $h$}{p11-Hfourier}
If the impulse response satisfies $h\in\ell^1$, the transfer function is the
(discrete) Fourier transform of $h$:
\[
H(f)=\sum_{m\in\Ints}h_m\,e^{-2\pi i f m},\qquad f\in\Reals .
\]
\textit{Proof idea:} feed the wave $\xi_{t,f}$ into the convolution
representation: $H(f)=\filt[\{\xi_{t,f}\}]_0=\sum_m h_{-m}e^{2\pi i f m}
=\sum_m h_m e^{-2\pi i f m}$ (re-index $m\mapsto -m$); $h\in\ell^1$ makes the sum
absolutely convergent. (Exercise.)\\
\textit{Why:} the practical recipe --- read off $h$ from the filter, then sum the
series to get $H$ (and conversely, recover $h$ as Fourier coefficients of $H$).
\end{theorem}

\begin{theorem}{Filtering preserves stationarity}{p11-statpres}
Let $\{X_t\}$ be stationary and $\filt$ an LTI filter with $h\in\ell^1$. Then
$\{Y_t\}=\filt[\{X_t\}]$ is stationary.\\[2pt]
\textit{Proof idea:} by the convolution theorem $Y_t\ms\sum_m h_m X_{t-m}$.
By Fubini, $\E[Y_t]=\sum_m h_m\E[X_{t-m}]=\mu_X\sum_m h_m$ is $t$-free; the
covariance
$\acvs^{(Y)}_\tau=\sum_{m,s}h_m\conj{h_s}\,\acvs^{(X)}_{\tau-m+s}$
depends only on $\tau$; and $\acvs^{(Y)}_0\le\acvs^{(X)}_0\norm{h}_1^2<\infty$,
so all is finite.\\
\textit{Why:} it licenses speaking of $\sdf_Y$ at all --- only a stationary
output has a spectral density.
\end{theorem}

\begin{theorem}{Output spectrum: $\sdf_Y(f)=\abs{H(f)}^2\sdf_X(f)$}{p11-outspec}
Let $\{X_t\}$ be stationary with $\acvs^{(X)}\in\ell^1$, and $\filt$ an LTI
filter with $h\in\ell^1$ and transfer function $H$. If
$\{Y_t\}=\filt[\{X_t\}]$ then
\[
\sdf_Y(f)=\abs{H(f)}^2\,\sdf_X(f),\qquad f\in\Reals .
\]
\textit{Proof idea:} from the covariance formula above,
$\acvs^{(Y)}_\tau=\sum_u w_u\,\acvs^{(X)}_{\tau-u}$ where
$w_u=\sum_s h_{u+s}\conj{h_s}$ (set $u=m-s$). With $\tilde h_t=\conj{h_{-t}}$ one
has $w=h*\tilde h\in\ell^1$, so the convolution theorem gives
$\sdf_Y=W\,\sdf_X$ with
$W(f)=H(f)\,\widetilde H(f)=H(f)\conj{H(f)}=\abs{H(f)}^2$, since
$\widetilde H(f)=\conj{H(f)}$.\\
\textit{Why:} the workhorse identity. It converts \emph{any} filtering question
into multiplying the input spectrum by a gain $\abs{H(f)}^2$ --- and, applied to
$\Phi(\Bsh)$ and $\Theta(\Bsh)$, it yields the ARMA spectral density below.
\end{theorem}

\begin{lemma}{Transfer function of a polynomial filter $\Phi(\Bsh)$}{p11-polyfilter}
For a degree-$p$ polynomial $\Phi(z)=\sum_{j=0}^p\phi_j z^j$, the filter
$\Phi(\Bsh)$ is LTI with transfer function
\[
H(f)=\Phi\!\big(e^{-2\pi i f}\big),\qquad f\in\Reals .
\]
\textit{Proof idea:} $\Phi(\Bsh)$ is a linear combination of the LTI filters
$\Bsh^j$, hence LTI, with impulse response $h_m=\phi_m$ for $0\le m\le p$ and $0$
otherwise. Then $H(f)=\sum_m h_m e^{-2\pi i f m}
=\sum_{j=0}^p\phi_j e^{-2\pi i f j}=\Phi(e^{-2\pi i f})$.\\
\textit{Why:} it turns the operator $\Phi(\Bsh)$ into a concrete function of
$f$; the bridge from time-domain difference equations to the spectrum.
\end{lemma}

\begin{theorem}{ARMA spectral density}{p11-arma}
For a stationary $\mathrm{ARMA}(p,q)$ process $\Phi(\Bsh)X_t=\Theta(\Bsh)\eps_t$
with $\eps_t$ white noise of variance $\sigma^2$,
\[
\sdf_X(f)=\sigma^2\,\frac{\abs{\Theta(e^{-2\pi i f})}^2}{\abs{\Phi(e^{-2\pi i f})}^2}
=\sigma^2\,\left|\frac{\sum_{j=0}^q\theta_j e^{-2\pi i j f}}
{\sum_{j=0}^p\phi_j e^{-2\pi i j f}}\right|^2 .
\]
\textit{Proof idea:} write $\Phi(\Bsh)[\{X_t\}]=\Theta(\Bsh)[\{\eps_t\}]$.
Applying the output-spectrum theorem to each side and the polynomial-filter
lemma,
$\abs{\Phi(e^{-2\pi i f})}^2\sdf_X(f)=\abs{\Theta(e^{-2\pi i f})}^2\sdf_\eps(f)$
with $\sdf_\eps(f)=\sigma^2$; divide, legal because stationarity forbids zeros
of $\Phi$ on the unit circle, so $\Phi(e^{-2\pi i f})\ne0$.\\
\textit{Why:} ARMA processes have \emph{no} closed-form ACVS but a clean
closed-form spectral density --- the practical way to read off their
frequency content.
\end{theorem}

% ----------------------------------------------------------------------
% ----------------------------------------------------------------------
